\chapter{Talking to the Source}

The last chapter ended with three ways to answer ``who talks to the plant?''.
This chapter takes the first and simplest one: every program that needs the
plant talks to it directly. Here that means two programs. The backend of the
operator screen reads the plant and sends it commands. The historian reads the
plant and keeps everything.

``Everyone asks the source'' had a bad reputation in the last chapter, but the
problem there was asking: polling, on a timer, for everything. Talking to the
source directly can be done well, and this chapter is about how.

\begin{figure}[htbp]
\centering
\begin{tikzpicture}
  \node[hf/source, minimum width=30mm, minimum height=11mm] (pl) {plant\\{\footnotesize OPC UA server}};
  \node[hf/reader, minimum width=30mm, right=34mm of pl, yshift=12mm] (bff) {screen backend};
  \node[hf/store, minimum width=30mm, right=34mm of pl, yshift=-14mm] (hi) {historian};
  \foreach \i/\dy in {1/9mm,2/0mm,3/-9mm} {
    \node[hf/reader, minimum width=18mm, minimum height=6mm, right=16mm of bff, yshift=\dy] (b\i) {browser};
    \draw[hf/push] (bff.east) -- (b\i.west);
  }
  \node[hf/reader, minimum width=30mm, right=16mm of hi] (sup) {supervisor};

  \draw[hf/push] (pl.north east) -- node[hf/note, above, sloped] {subscription} ([yshift=1.5mm]bff.west);
  \draw[hf/ask] ([yshift=-1.5mm]bff.west) -- node[hf/note, below, sloped] {commands} ([yshift=-2.5mm]pl.north east);
  \draw[hf/push] (pl.south east) -- node[hf/note, below, sloped] {subscription} (hi.west);
  \draw[hf/ask] (sup) -- (hi);
\end{tikzpicture}
\caption{Direct connection. Two programs hold a session with the plant. The
browsers never reach the plant: the screen backend serves all of them.}
\end{figure}

\section{The simplest thing that works}

A broker in the middle is the right answer when there are many publishers, many
readers, or links that drop. A simulated plant running in the same cluster as
two readers has none of those problems. Putting a broker there means running,
monitoring and securing a piece of infrastructure to solve a problem that
doesn't exist.

\begin{history}{1998--1999 \textperiodcentered{} Extreme Programming}
In the late 1990s, Kent Beck, Ron Jeffries and others were building Extreme
Programming, a way of working that rejected big designs made up front. In 1998,
Jeffries wrote a short essay called ``You're NOT gonna need it!''. Beck's
\emph{Extreme Programming Explained}, in 1999, gave the practice its acronym:
\textbf{YAGNI}, \emph{you aren't going to need it}.
\end{history}

\begin{concept}{YAGNI}
Build what the problem in front of you needs, not what you predict a future
problem will need.

\textbf{What it fixes:} you stop paying, every day, for infrastructure that
answers questions nobody is asking.

\textbf{What it costs:} when the future problem does arrive, you pay to add the
solution then. That cost stays small if the contracts between parts stay the
same.
\end{concept}

The last part is what makes YAGNI safe here. The readers depend on the plant's
OPC UA address space: its nodes, their names and types. If a third reader shows
up one day, or the load of the sessions becomes measurable, a broker can be put
in front of that same address space without changing what the readers expect.

\section{What a connection is made of}

An OPC UA connection isn't one thing. It's four layers, each with its own life.

\begin{figure}[htbp]
\centering
\begin{tikzpicture}
  \node[draw=black!40, line width=0.8pt, rounded corners=4pt, minimum width=124mm, minimum height=48mm, label={[hf/note, anchor=north west]north west:TCP connection}] (tcp) {};
  \node[draw=hfSource, line width=0.8pt, rounded corners=4pt, minimum width=114mm, minimum height=38mm, label={[hf/note, anchor=north west, text=hfSource]north west:secure channel: encryption, renewed periodically}] at ([yshift=-2.5mm]tcp.center) (sc) {};
  \node[draw=hfReader, fill=hfReader!6, line width=0.8pt, rounded corners=4pt, minimum width=104mm, minimum height=28mm, label={[hf/note, anchor=north west, text=hfReader]north west:session: who the client is, survives a lost channel}] at ([yshift=-2.5mm]sc.center) (se) {};
  \node[draw=hfBus, fill=hfBus!8, line width=0.8pt, rounded corners=4pt, minimum width=44mm, minimum height=14mm, label={[hf/note, anchor=north west, text=hfBus]north west:subscription}] at ([xshift=-25mm, yshift=-3.5mm]se.center) (s1) {};
  \node[draw=hfBus, fill=hfBus!8, line width=0.8pt, rounded corners=4pt, minimum width=44mm, minimum height=14mm, label={[hf/note, anchor=north west, text=hfBus]north west:subscription}] at ([xshift=25mm, yshift=-3.5mm]se.center) (s2) {};
  \foreach \s in {s1,s2} \foreach \dx in {-12mm,0mm,12mm}
    \node[draw=hfBus, fill=white, minimum width=9mm, minimum height=4mm, inner sep=1pt, font=\sffamily\tiny] at ([xshift=\dx, yshift=-2mm]\s.center) {item};
\end{tikzpicture}
\caption{The layers of an OPC UA connection. Monitored items live inside
subscriptions, subscriptions inside a session, the session on a secure channel,
and the channel on a TCP connection.}
\end{figure}

\begin{concept}{Session}
A session identifies a client to the server: who it is and what it may do. It
lives on a secure channel but doesn't die with it. If the network drops, the
client opens a new channel and \emph{reactivates} the same session, as long as
it does so before the session timeout. Its subscriptions are still there,
waiting.

\textbf{What it fixes:} a short network glitch doesn't throw away the client's
state, so nothing has to be set up again.

\textbf{What it costs:} the server keeps that state in memory for every client,
for as long as the timeout allows.
\end{concept}

This is the part most clients get wrong. A client that treats every error as
``start over'' creates a new session each time, leaves the old one hanging on
the server until it times out, and loses whatever changed in between.

\section{Subscribe, don't poll}

The previous chapter showed subscriptions from the outside. Here are the knobs.
Each monitored item has a \emph{sampling interval} (how often the server looks
at the value), a \emph{deadband} (how much it must change to count) and a
\emph{queue size} (how many changes to keep between two deliveries). Each
subscription has a \emph{publishing interval} (how often the server sends what
it has).

The two readers want different things, so they set the knobs differently.

\begin{itemize}
  \item \textbf{The screen wants the present.} A queue of one is right: if a
        value changed twice between deliveries, the screen only needs the
        latest. The deadband can be generous, at the level a person would
        notice.
  \item \textbf{The historian wants everything.} Its queue must hold every
        sample that can happen between two deliveries, or changes are thrown
        away before they arrive. Its deadband is the exception reporting from
        the last chapter, set close to the instrument's accuracy.
\end{itemize}

There's one more detail that matters for a screen. When a monitored item is
created, the server always reports its current value first, whether or not it
changes afterwards. A screen that connects while the plant is steady gets the
whole picture immediately, including the analyzer that only updates every six
minutes.

\section{Telling a steady plant from a dead server}

When nothing changes, a subscription sends nothing. That raises the same
question every push system has to answer: is the plant steady, or is the
server gone?

\begin{concept}{Keep-alive and lifetime}
If a subscription has had nothing to send for a number of publishing intervals
(the \emph{max keep-alive count}), the server sends an empty message anyway, a
\emph{keep-alive}. The client knows that, so when even the keep-alive is late,
the server is gone.

It works the other way too. If the server gets no requests from the client for
a number of intervals (the \emph{lifetime count}), it deletes the subscription.
The specification requires the lifetime to be at least three times the
keep-alive, so the client always has time to notice before its subscription
disappears.

\textbf{What it fixes:} silence stops being ambiguous. Each side knows, within a
bounded time, that the other is gone.

\textbf{What it costs:} a little traffic when nothing is happening, and two
more numbers to tune.
\end{concept}

\begin{figure}[htbp]
\centering
\begin{tikzpicture}[x=0.95cm]
  \draw[black!40, -{Stealth[length=2mm]}] (0,0) -- (14.6,0) node[hf/note, below left] {time};
  \foreach \t in {1,...,14} \draw[black!40] (\t,0.08) -- (\t,-0.08);
  \node[hf/note, anchor=north] at (1,-0.15) {publishing\\interval};
  % data
  \foreach \t in {1,2,4} {
    \fill[hfBus] (\t,0.55) circle (2.4pt);
  }
  \node[hf/note, text=hfBus, anchor=south] at (2.5,0.8) {changes};
  % keep-alive after 3 empty intervals
  \draw[hfBus, fill=white, line width=0.8pt] (7,0.55) circle (2.4pt);
  \node[hf/note, anchor=south east] at (7.2,0.8) {keep-alive};
  \node[hf/note, anchor=north] at (6,-0.15) {3 empty};
  % server dies
  \draw[hf/pain] (8.25,0.3) -- (8.75,0.8) (8.25,0.8) -- (8.75,0.3);
  \node[hf/note, text=hfAlert, anchor=south west] at (8.2,0.9) {server gone};
  % expected keep-alive missing
  \draw[hfAlert, densely dashed, line width=0.9pt] (10,0.55) circle (4pt);
  \node[hf/note, text=hfAlert, anchor=north] at (10.6,-0.15) {keep-alive late:\\client marks data bad};
\end{tikzpicture}
\caption{Keep-alive with a count of three. After three empty intervals the
server says ``still here''. When that message doesn't come, the client knows.}
\end{figure}

What the client does next is the important part.

\begin{concept}{Data quality}
Since OPC DA, every industrial value has traveled with a \emph{quality}: good,
bad or uncertain. In OPC UA every value comes as a \emph{data value}: the value,
a status code, the time it was measured at the source and the time the server
handled it. A screen that shows a number is claiming that the plant is doing
that \emph{now}, and when the source stops, the screen has to stop making the
claim.

\textbf{What it fixes:} the operator can tell a steady plant from a silent one.

\textbf{What it costs:} every layer between the plant and the screen has to
carry the quality along, not just the value.
\end{concept}

When the connection is lost, the screen backend marks every value as bad, and
the screen shows ``no data''. A frozen value on an operator screen is worse than
an empty one, because it looks exactly like a plant that is holding steady.

For the historian the rule is the same, applied to history: a period without
data is stored as a gap, marked bad. Drawing a line across it would invent
values the plant never had.

\section{When the connection drops}

Giving up after a few tries is the worst policy for a program whose job is to
watch a plant. The plant will come back, and the screen has to be there when it
does. So the client retries for as long as it runs. The question is how fast.

\begin{history}{1970--2015 \textperiodcentered{} Backing off}
In 1970, ALOHAnet linked computers across the Hawaiian islands by radio, and
when two stations collided, each waited a random time before sending again. In
1976, Robert Metcalfe and David Boggs described Ethernet in \emph{Communications
of the ACM}, and Ethernet made the waiting smarter: after each new collision, a
station doubles the range of its random wait. That's \textbf{binary exponential
backoff}. In 2015, Marc Brooker wrote ``Exponential Backoff and Jitter'' on the
AWS Architecture Blog, with simulations showing that the randomness, the
\emph{jitter}, matters as much as the doubling.
\end{history}

\begin{concept}{Exponential backoff with jitter}
After each failed attempt, wait longer than the last time, usually twice as
long, up to a ceiling. Then pick a random wait below that value instead of
using it exactly.

\textbf{What it fixes:} a client that's down doesn't hammer a server that's
trying to come back, and many clients that failed together don't all retry at
the same instant.

\textbf{What it costs:} after a long outage, recovery takes up to one ceiling
longer than it strictly had to.
\end{concept}

\begin{figure}[htbp]
\centering
\begin{tikzpicture}[x=0.42cm, y=0.5cm]
  \draw[black!40, -{Stealth[length=2mm]}] (0,0) -- (32,0);
  \node[hf/note, anchor=north] at (16,-0.9) {seconds after the failure};
  \foreach \t in {0,5,...,30} \draw[black!40] (\t,0.1) -- (\t,-0.1) node[hf/note, below] {\t};
  % without jitter: 1,2,4,8,16 cumulative 1,3,7,15,31 -> show 1,3,7,15,(31 off)
  \node[hf/note, anchor=east] at (-0.5,3) {no jitter};
  \foreach \t in {1,3,7,15,30} \foreach \k in {0,1,2,3,4} \fill[hfAlert] (\t,2.4+0.3*\k) circle (2.2pt);
  \node[hf/note, anchor=east] at (-0.5,1.1) {with jitter};
  \foreach \t/\y in {0.4/0.6, 0.9/1.2, 0.7/1.8, 1.8/0.9, 2.6/1.5, 3.4/0.6, 4.9/1.8, 5.8/1.2, 2.1/1.8, 6.6/0.9,
                     9.2/1.5, 11.5/0.6, 8.1/1.2, 12.9/1.8, 14.2/0.9, 17.0/1.5, 21.4/0.6, 19.6/1.2, 25.3/1.8, 27.9/0.9} {
    \fill[hfReader] (\t,\y) circle (2.2pt);
  }
\end{tikzpicture}
\caption{Five clients lose the same server. Without jitter they retry at the
same instants, 1, 3, 7, 15 and 30 seconds, and hit the server together every
time. With jitter the attempts spread out.}
\end{figure}

Reconnecting well means doing it in order. First, try to reactivate the session
on a new channel. If the session is still alive, the subscriptions are too, and
the client asks the server to \emph{republish} the notifications it missed:
every notification carries a sequence number, so the client knows exactly which
ones didn't arrive. Only when the session is gone does the client create a new
one and rebuild its subscriptions. For the historian, this order is the
difference between a complete history and one with holes in it.

\section{A backend for the screen}

The browsers don't talk to the plant. Between them sits a backend made for that
one screen.

\begin{history}{2013--2015 \textperiodcentered{} SoundCloud}
SoundCloud had one general-purpose API serving its website, its iPhone app and
its Android app. Each client needed the data in a different shape, every screen
change waited on the API team, and the API grew into a compromise that fit
nobody. Around 2013, the company started giving each frontend its own backend,
owned by the team that built that frontend. Phil Calçado described the pattern
in 2015 as the \textbf{Backend for Frontend}, or BFF, and Sam Newman's write-up
the same year spread it.
\end{history}

\begin{concept}{Backend for Frontend (BFF)}
A backend built for exactly one kind of client. It talks to whatever services
exist behind it and hands that client the data in the shape the screen needs,
over the transport the screen can use.

\textbf{What it fixes:} the screen stops depending on contracts designed for
somebody else, and internal contracts stay internal.

\textbf{What it costs:} one more service to run, and the risk that logic which
belongs to the domain leaks into it because it's the easiest place to put it.
\end{concept}

Here the BFF has one more job: it keeps the number of sessions fixed. Ten
operators with ten browsers still mean one session against the plant, because
the BFF holds that session and fans the data out itself. It keeps the last data
value of every node in memory. When a browser connects, it sends that whole
picture at once, and from then on only changes. When the connection to the
plant is lost, it tells every browser at once.

How the BFF pushes to the browser, WebSocket or Server-Sent Events, is a story
of its own.

\section{Commands}

The screen also needs to act on the plant: pause it, change its speed, switch a
disturbance on. Those commands go through the same session, but they are a
different kind of conversation.

\begin{history}{1988 \textperiodcentered{} Bertrand Meyer}
In \emph{Object-Oriented Software Construction}, Bertrand Meyer stated the
\textbf{command--query separation} principle: a method either changes the state
of an object or answers a question about it, never both. Around 2010, Greg
Young took the idea from methods up to whole systems and called it CQRS: one
path for commands, another for queries.
\end{history}

A subscription is a query that never ends. A command is a single request that
needs an answer: did the plant accept it? OPC UA has two ways to send one. A
\emph{Write} sets the value of a node and returns a status code for it. A
\emph{Call} invokes a \emph{method} on the server, with arguments and a result,
like a function.

Use a write for a value, like a setpoint or a speed. Use a method for an
action, like ``pause'' or ``start disturbance 6''. An action written as a
variable (``set \texttt{paused} to true'') looks simple until two clients write
it at the same time, or until someone needs the action to refuse with a reason.

\begin{figure}[htbp]
\centering
\begin{tikzpicture}[x=1cm, y=-0.72cm]
  \node[hf/reader, minimum width=28mm] at (0,0) {screen backend};
  \node[hf/source, minimum width=24mm] at (7,0) {plant};
  \draw[hf/lifeline] (0,0.7) -- (0,6.4);
  \draw[hf/lifeline] (7,0.7) -- (7,6.4);
  \draw[hf/push] (7,1.4) -- node[hf/note, above] {notifications} (0,1.4);
  \draw[hf/push] (7,2.2) -- (0,2.2);
  \draw[hf/ask] (0,3.3) -- node[hf/note, above] {Call \texttt{Pause()}} (7,3.3);
  \draw[hf/answer] (7,4.1) -- node[hf/note, above] {status: Good} (0,4.1);
  \draw[hf/push] (7,5.0) -- node[hf/note, above] {notification: paused} (0,5.0);
  \draw[hf/push] (7,5.8) -- (0,5.8);
\end{tikzpicture}
\caption{One session, two kinds of conversation. Notifications flow without
being asked. A command is a request that gets its own answer.}
\end{figure}

\section{What it costs}

The plant holds one session per reader, and each reader has its own way of
failing. With two readers, each doing reconnection properly, that's a cost
worth paying for a system with fewer moving parts. A third reader means a third
session, and at some point the count, or the load it puts on the plant, becomes
a real problem. That's when the bus from the last chapter starts paying for
itself.

\section{Writing the choice down}

Choosing this architecture over the other two is a decision someone will
question later, and the costs above are exactly what they'll want to ``fix''.
There's a standard way to keep the reasons around.

\begin{history}{2011 \textperiodcentered{} Michael Nygard}
In November 2011, Michael Nygard, then at Relevance, published a short blog
post called ``Documenting Architecture Decisions''. His complaint was about what
happens to a project after the people who designed it leave. A new developer
finds a strange choice and has two options, both bad: accept it blindly,
because there might be a reason, or reverse it blindly, because nobody
remembers one. His answer was the \textbf{Architecture Decision Record}.
\end{history}

\begin{concept}{Architecture Decision Record (ADR)}
A short, numbered text file kept in the repository, one per decision, with a
fixed shape: the \emph{context} (what forced the decision), the
\emph{decision}, its \emph{status} (proposed, accepted, superseded) and its
\emph{consequences}, good and bad. An ADR is never edited to change its mind. A
new ADR supersedes it.

\textbf{What it fixes:} the reason behind a choice survives the people who made
it.

\textbf{What it costs:} you have to write while deciding, when you'd rather be
coding.
\end{concept}

The most useful part is the bad consequences. Writing them down shows that a
cost was accepted on purpose, so whoever comes later knows it isn't a bug.

Reversals belong in the record too. When a team adopts a broker and then
decides it doesn't need one, the first ADR stays, marked superseded, and the
second one says why. Whoever proposes the broker again will find both, along
with the reason it was dropped.

\section*{Further reading}

\begin{itemize}
  \item OPC Foundation, \emph{OPC UA Part 4: Services}, sections on sessions,
        subscriptions and monitored items,
        \url{https://reference.opcfoundation.org/Core/Part4/}.
  \item K.\ Beck, \emph{Extreme Programming Explained}, Addison-Wesley, 1999.
  \item R.~M.\ Metcalfe and D.~R.\ Boggs, ``Ethernet: Distributed Packet
        Switching for Local Computer Networks'', \emph{Communications of the
        ACM} 19(7), 1976.
  \item M.\ Brooker, ``Exponential Backoff and Jitter'', AWS Architecture Blog,
        2015,
        \url{https://aws.amazon.com/blogs/architecture/exponential-backoff-and-jitter/}.
  \item P.\ Calçado, ``The Back-end for Front-end Pattern (BFF)'', 2015,
        \url{https://philcalcado.com/2015/09/18/the_back_end_for_front_end_pattern_bff.html}.
  \item B.\ Meyer, \emph{Object-Oriented Software Construction}, Prentice Hall,
        1988.
  \item M.\ Nygard, ``Documenting Architecture Decisions'', 2011,
        \url{https://cognitect.com/blog/2011/11/15/documenting-architecture-decisions}.
\end{itemize}
